\section*{曲率与曲率圆}

1. 求下列曲线在指定点处的曲率。

(1) 曲线 \(xy=4\)，在点 \((2,2)\) 处。

解：由 \(xy=4\) 得 \(y=\frac{4}{x}\)，求导：
\[ y'=-\frac{4}{x^{2}},\qquad y''=\frac{8}{x^{3}}. \]
在点 \((2,2)\) 处：
\[ y'(2)=-\frac{4}{4}=-1,\qquad y''(2)=\frac{8}{8}=1. \]
曲率公式 \(K=\frac{|y''|}{(1+y'^{2})^{3/2}}\)，故
\[ K=\frac{1}{(1+(-1)^{2})^{3/2}}=\frac{1}{2^{3/2}}=\frac{1}{2\sqrt{2}}=\frac{\sqrt{2}}{4}. \]

(2) 曲线 \(y=4x-x^{2}\)，在点 \((0,0)\) 处。

解：求导：
\[ y'=4-2x,\qquad y''=-2. \]
在点 \((0,0)\) 处：
\[ y'(0)=4,\qquad y''(0)=-2. \]
曲率：
\[ K=\frac{|-2|}{(1+4^{2})^{3/2}}=\frac{2}{17^{3/2}}=\frac{2}{17\sqrt{17}}. \]

(3) 曲线 \(y=\ln\left(x+\sqrt{1+x^{2}}\right)\)，在点 \((0,0)\) 处。

解：求导：
\[ y'=\frac{1+\frac{x}{\sqrt{1+x^{2}}}}{x+\sqrt{1+x^{2}}}=\frac{1}{\sqrt{1+x^{2}}},\qquad y''=-\frac{x}{(1+x^{2})^{3/2}}. \]
在点 \((0,0)\) 处：
\[ y'(0)=1,\qquad y''(0)=0. \]
曲率：
\[ K=\frac{|0|}{(1+1^{2})^{3/2}}=0. \]

(4) 曲线 \(y=\ln x\)，在点 \((1,0)\) 处。

解：求导：
\[ y'=\frac{1}{x},\qquad y''=-\frac{1}{x^{2}}. \]
在点 \((1,0)\) 处：
\[ y'(1)=1,\qquad y''(1)=-1. \]
曲率：
\[ K=\frac{|-1|}{(1+1^{2})^{3/2}}=\frac{1}{2^{3/2}}=\frac{1}{2\sqrt{2}}=\frac{\sqrt{2}}{4}. \]

2. 求由下列参数方程表示的曲线在指定参数处的曲率。

(1) 曲线
\[ \begin{cases} x=3t^{2},\\ y=3t-t^{3}, \end{cases}\qquad t=1. \]

解：求导：
\[ x'=6t,\quad x''=6;\qquad y'=3-3t^{2},\quad y''=-6t. \]
在 \(t=1\) 处：
\[ x'(1)=6,\quad x''(1)=6;\qquad y'(1)=3-3=0,\quad y''(1)=-6. \]
参数方程的曲率公式 \(K=\frac{|x'y''-x''y'|}{(x'^{2}+y'^{2})^{3/2}}\)，故
\[ K=\frac{|6\cdot(-6)-6\cdot 0|}{(6^{2}+0^{2})^{3/2}}=\frac{36}{36^{3/2}}=\frac{36}{216}=\frac{1}{6}. \]

(2) 曲线
\[ \begin{cases} x=a(\cos t+\sin t),\\ y=a(\sin t-t\cos t), \end{cases}\qquad t=\frac{\pi}{2}\ (a>0). \]

解：求导：
\[ x'=a(-\sin t+\cos t),\qquad x''=a(-\cos t-\sin t); \]
\[ y'=a(\cos t-(\cos t-t\sin t))=at\sin t,\qquad y''=a(\sin t+t\cos t). \]
在 \(t=\frac{\pi}{2}\) 处：
\[ x'\left(\frac{\pi}{2}\right)=a(-1+0)=-a,\qquad x''\left(\frac{\pi}{2}\right)=a(0-1)=-a; \]
\[ y'\left(\frac{\pi}{2}\right)=a\cdot\frac{\pi}{2}\cdot 1=\frac{a\pi}{2},\qquad y''\left(\frac{\pi}{2}\right)=a(1+0)=a. \]
曲率：
\[ \begin{aligned} K&=\frac{|x'y''-x''y'|}{(x'^{2}+y'^{2})^{3/2}}
=\frac{\left|(-a)\cdot a-(-a)\cdot\frac{a\pi}{2}\right|}{\left((-a)^{2}+\left(\frac{a\pi}{2}\right)^{2}\right)^{3/2}}\\
&=\frac{a^{2}\left|-1+\frac{\pi}{2}\right|}{\left(a^{2}+\frac{a^{2}\pi^{2}}{4}\right)^{3/2}}
=\frac{\left|\frac{\pi}{2}-1\right|}{a\left(1+\frac{\pi^{2}}{4}\right)^{3/2}}. \end{aligned} \]

4. 求曲线 \(y=x^{2}\) 上任一点处的曲率，并问哪一点处曲率最大？

解：求导：
\[ y'=2x,\qquad y''=2. \]
曲率：
\[ K=\frac{|y''|}{(1+y'^{2})^{3/2}}=\frac{2}{(1+4x^{2})^{3/2}}. \]
为求曲率最大点，令 \(u=1+4x^{2}\)，则 \(K=2u^{-3/2}\)。对 \(K\) 关于 \(x\) 求导：
\[ \frac{dK}{dx}=2\cdot\left(-\frac{3}{2}\right)u^{-5/2}\cdot 8x=-\frac{24x}{u^{5/2}}. \]
令 \(\frac{dK}{dx}=0\)，得 \(x=0\)。当 \(x<0\) 时 \(\frac{dK}{dx}>0\)；当 \(x>0\) 时 \(\frac{dK}{dx}<0\)。故在 \(x=0\) 处曲率最大，最大曲率
\[ K_{\max}=\frac{2}{(1+0)^{3/2}}=2. \]

5. 求椭圆周 \(\frac{x^{2}}{a^{2}}+\frac{y^{2}}{b^{2}}=1\) 上任一点处的曲率，并问哪一点处曲率最大？其中 \(a>b>0\)。

解：将椭圆表示为参数方程：
\[ x=a\cos t,\qquad y=b\sin t. \]
求导：
\[ x'=-a\sin t,\quad x''=-a\cos t;\qquad y'=b\cos t,\quad y''=-b\sin t. \]
参数方程的曲率：
\[ \begin{aligned} K&=\frac{|x'y''-x''y'|}{(x'^{2}+y'^{2})^{3/2}}
=\frac{|(-a\sin t)(-b\sin t)-(-a\cos t)(b\cos t)|}{(a^{2}\sin^{2}t+b^{2}\cos^{2}t)^{3/2}}\\
&=\frac{ab(\sin^{2}t+\cos^{2}t)}{(a^{2}\sin^{2}t+b^{2}\cos^{2}t)^{3/2}}
=\frac{ab}{(a^{2}\sin^{2}t+b^{2}\cos^{2}t)^{3/2}}. \end{aligned} \]
为求曲率最大点，需使分母最小，即求 \(f(t)=a^{2}\sin^{2}t+b^{2}\cos^{2}t\) 的最小值。令 \(u=\sin^{2}t\)，则
\[ f(t)=a^{2}u+b^{2}(1-u)=b^{2}+(a^{2}-b^{2})u. \]
由于 \(a>b>0\)，\(a^{2}-b^{2}>0\)，故当 \(u=0\) 即 \(\sin t=0\) 时，\(f(t)\) 最小。此时 \(t=0\) 或 \(\pi\)，对应点 \((a,0)\) 和 \((-a,0)\)。最大曲率：
\[ K_{\max}=\frac{ab}{(b^{2})^{3/2}}=\frac{ab}{b^{3}}=\frac{a}{b^{2}}. \]
故在椭圆长轴端点处曲率最大。

7. 求曲线 \(y=\ln x\) 与 \(x\) 轴交点处的曲率圆方程。

解：曲线与 \(x\) 轴交点：令 \(y=0\)，则 \(\ln x=0\)，\(x=1\)，交点为 \((1,0)\)。

由第 1 题 (4) 知，在点 \((1,0)\) 处曲率 \(K=\frac{\sqrt{2}}{4}\)。曲率半径
\[ \rho=\frac{1}{K}=\frac{4}{\sqrt{2}}=2\sqrt{2}. \]
在点 \((1,0)\) 处，切线斜率 \(y'=\frac{1}{x}=1\)，法线斜率为 \(-1\)，法线方程
\[ y-0=-1(x-1), \]
即 \(y=-x+1\)。
曲率圆心在法线上，且与切点的距离为 \(\rho=2\sqrt{2}\)，故圆心坐标有两种可能：
\[ (x_0,y_0)=(1,0)\pm\frac{2\sqrt{2}}{\sqrt{1^{2}+(-1)^{2}}}(-1,1)=(1,0)\pm 2(-1,1), \]
即 \((1\mp 2,\pm 2)\)，也就是 \((-1,2)\) 或 \((3,-2)\)。

由于 \(y''=-\frac{1}{x^{2}}=-1<0\)，曲线凹向下，故圆心应位于曲线下方，取 \((3,-2)\)。曲率圆方程为
\[ (x-3)^{2}+(y+2)^{2}=(2\sqrt{2})^{2}=8. \]
